\chapter{Perturbed multistart batch (menu 29)}
\label{ch:menu29}
\menulabel{29}

\section{Purpose}

Hartree--Fock and Kohn--Sham SCF allow multiple self-consistent solutions;
converging to a non-global minimum delivers wrong energies and properties and
is easy to miss.  The source's inexpensive test (Vaucher \& Reiher, JCTC
2017, \S IV) perturbs a converged solution's orbitals --- random
occupied--virtual pair mixings, Eqs.~(2)--(3), 10 random pairs out of the 15
highest occupied and the 15 lowest unoccupied orbitals with random angles in
$[0, 90)$ degrees --- and re-runs: a perturbed start either returns to the
same solution or finds another stationary point, and a lower energy
identifies the reference as wrongly converged.  Stability analysis detects
unstable solutions (saddle points) but cannot distinguish local from global
minima, so this is complementary.  The menu generates the perturbed starts
(the engine run stays with the user); the perturbation is fully recorded
(seed, every pair, every angle) and replays byte-identically.

\section{What you need}

The converged reference's Molekel mkl (export it with
\code{orca\_2mkl <base> -mkl}) and the ORCA input that produced it (the menu
rewrites that input's guess; its inline geometry is cross-checked against the
reference's coordinates).  The number of starts, the random seed, the pairs
per start and the window per side are menu options (defaults 3, 20260927,
10 and 15 --- the last two the source's values).

\section{Walkthrough}

\lstinputlisting[style=fbkcode, firstline=34, lastline=68,
  title={The menu-29 example session (the CH$_4$ dissociation reference: the
  UKS propagated guess; three perturbed starts are generated, each with its
  windows, pairs and angles recorded)}]
  {examples/expected/m29_perturb.txt}

\section{What you get}

For every start: \file{<reference>.p<k>.fbk.mkl} (the perturbed orbitals,
gbw-ready) and \file{<input>.p<k>.inp} (the base input with \code{!MORead}
and \code{\%moinp "<reference>.p<k>.fbk.gbw"} in place), plus a report
\file{<input>.perturb.fbk.md} with the full record.  Next steps: convert each
orbital file next to itself with \code{orca\_2mkl <name>.fbk -gbw}, run the
inputs, and compare the converged energies with the reference's --- a lower
energy heals a wrongly converged reference, the same (or a higher) energy is
no information.

\section{Options worth knowing}

\begin{readbox}
\begin{itemize}
  \item \textbf{measured anchor} (CH$_4$ dissociation, UKS PBE0/def2-SVP
    with one C--H at 2.6~\AA, fixtures \file{ch4\_*}): the guess propagated
    from the equilibrium orbitals keeps the restricted solution
    ($-40.111230225721$~Eh, $\langle S^2\rangle = 0.000000$); ORCA's default
    guess also lands on a restricted solution
    ($-40.183584827774$~Eh, $\langle S^2\rangle = 0.000000$); all three
    generated starts reach the broken-symmetry solution
    ($-40.2416$~Eh, $\langle S^2\rangle = 0.971943$) --- a healing of
    0.130~Eh and a further 0.058~Eh below the default guess;
  \item \textbf{the seed is the whole story}: each start draws from its own
    stream (\code{seed + k - 1}), every pair and angle is printed and
    written into the report, and the same seed reproduces the same files
    byte for byte;
  \item \textbf{both spins are perturbed} for an unrestricted reference
    (the alpha and beta blocks are mixed independently); a restricted
    reference has only the alpha block;
  \item \textbf{boundaries}: the test cannot guarantee detection (the
    source's own statement --- it is a verification, not a proof); the
    perturbed columns are mixtures, not eigenfunctions, so the occupations
    and energies in the file are the reference's and ORCA re-determines
    everything after reading the guess; the source validated the mechanism
    at the SCF level and notes that MC-SCF solutions can in principle be
    caught in local minima as well.
\end{itemize}
\end{readbox}
